\documentclass[12pt]{article}
\usepackage[margin=1in]{geometry}
\usepackage{tikz}
\usetikzlibrary{arrows.meta}
\usepackage[american voltages, european currents]{circuitikz}
\usepackage{amsmath}
\usepackage{amssymb}
\usepackage[table,dvipsnames]{xcolor}
\usepackage{pgfplots}
\pgfplotsset{compat=1.18}
\usepackage{float}
\usepackage{caption}
\usepackage{parskip}

% -- colors from source file --
\definecolor{myblue}{RGB}{0,103,172}
\definecolor{mygray}{RGB}{240,240,240}

\begin{document}

\begin{center}
  {\LARGE\textbf{TikZ / CircuitiKZ Figures}}\\[0.4em]
  {\large\texttt{tikz\_figures.tex}}
\end{center}
\vspace{1em}
\noindent Edit the TikZ code in this file, then recompile
with \texttt{pdflatex} to preview changes.  Run
\texttt{extract\_tikz\_figures.py} to regenerate the PNGs.

\clearpage

\noindent\textbf{Label:} \texttt{fig:unity-gain}\\[0.5em]
\begin{figure}[H]
  \centering
  \begin{circuitikz}[american, scale=1]
      			\ctikzset{bipoles/oscope/waveform=none}
      			\draw (-1,0) coordinate(NonInv) to[short, o-] ++ (2,0) node[above] {$3$}
      			node[op amp, noinv input up, anchor=+](OA){OA}
      			(OA.out) node[above] {$1$} to[short, -*] ++(1.5,0) coordinate(Out) to[short, -o] ++(1,0) node[right]{$\text{V}_{\text{out}}$}
      			(OA.-) node[above] {$2$} to[short,] ++(0,-1) to[short,] ++(2.75,0) to[short, -*] ++(0,1.5);
      			\draw (-3,-2) node[ground]{} to[sinusoidal voltage source] ++ (0,+2) to[short, -*] ++(2.0,0) coordinate(In1)
      			to[oscope=$\text{CH2}$, -] ++(0,-2) node[ground]{}
      			(Out) to[oscope=$\text{CH1}$, -] ++(0,-1.5) node[ground]{}
      			(In1) node[above] {$\text{V}_{\text{in}}$};
      		\end{circuitikz}
  \caption{Unity-gain buffer (voltage follower) circuit}
\end{figure}

\clearpage

\noindent\textbf{Label:} \texttt{fig:buffer-example}\\[0.5em]
\begin{figure}[H]
  \centering
  \begin{circuitikz}[american, scale=1]
      			\ctikzset{bipoles/oscope/waveform=none}
      			\draw (0,0) coordinate(NonInv) to[short, *-] ++ (0.5,0) node[above] {$3$}
      			node[op amp, noinv input down, anchor=+](OA){OA}
      			(OA.out) node[above] {$1$} to[short, -*] ++(1.5,0) coordinate(Out1) to[short, -*] ++(2,0) coordinate(Out2) to[short, -o] ++(0.5,0) node[right]{$\text{V}_{\text{out}}$}
      			(OA.-) node[below] {$2$} to[short,] ++(0,+1) to[short,] ++(2.75,0) to[short, -*] ++(0,-1.5);
      			\draw (-4,-2) node[ground]{} to[sinusoidal voltage source] ++ (0,+2)
      			to[R=$4.7 \text{ k}\Omega$, -*] ++(2,0)
      			to[R=$4.7 \text{ k}\Omega$] ++(0,-2) node[ground]{}
      			(Out2) to[oscope=$\text{CH1}$, -] ++(0,-2.5) node[ground]{}
      			(Out1) to[R=$10 \text{ k}\Omega$] ++(0,-2.5) node[ground]{}
      			(NonInv) to[oscope=$\text{CH2}$, -] ++(0,-2) node[ground]{}
      			(NonInv) to[short,] ++(-2,0)
      			(NonInv) node[above] {$\text{V}_{\text{in}}$};
      		\end{circuitikz}
  \caption{Buffer Circuit \#1 with the 10~k$\Omega$ load resistor.}
\end{figure}

\clearpage

\noindent\textbf{Label:} \texttt{fig:buffer-example2}\\[0.5em]
\begin{figure}[H]
  \centering
  \begin{circuitikz}[american, scale=1]
      				\ctikzset{bipoles/oscope/waveform=none}
      				\draw (0,0) coordinate(NonInv) to[short, *-] ++ (0.5,0) node[above] {$3$}
      				node[op amp, noinv input down, anchor=+](OA){OA}
      				(OA.out) node[above] {$1$} to[short, -*] ++(1.5,0) coordinate(Out1) to[short, -*] ++(2,0) coordinate(Out2) to[short, -o] ++(0.5,0) node[right]{$\text{V}_{\text{out}}$}
      				(OA.-) node[below] {$2$} to[short,] ++(0,+1) to[short,] ++(2.75,0) to[short, -*] ++(0,-1.5);
      				\draw (-4,-2) node[ground]{} to[sinusoidal voltage source] ++ (0,+2)
      				to[R=$4.7 \text{ k}\Omega$, -*] ++(2,0)
      				to[R=$4.7 \text{ k}\Omega$] ++(0,-2) node[ground]{}
      				(Out2) to[oscope=$\text{CH1}$, -] ++(0,-2.5) node[ground]{}
      				(Out1) to[R=$1 \text{ k}\Omega$] ++(0,-2.5) node[ground]{}
      				(NonInv) to[oscope=$\text{CH2}$, -] ++(0,-2) node[ground]{}
      				(NonInv) to[short,] ++(-2,0)
      				(NonInv) node[above] {$\text{V}_{\text{in}}$};
      			\end{circuitikz}
  \caption{Buffer Circuit \#2 with the 1~k$\Omega$ load resistor.}
\end{figure}

\clearpage

\noindent\textbf{Label:} \texttt{fig:buffer-example3}\\[0.5em]
\begin{figure}[H]
  \centering
  \begin{circuitikz}[american, scale=1]
      				\ctikzset{bipoles/oscope/waveform=none}
      				\draw (1.5,0) coordinate(NonInv) to[short, *-] ++ (0.5,0) node[above] {$3$}
      				node[op amp, noinv input down, anchor=+](OA){OA}
      				(OA.out) node[above] {$1$} to[short, -*] ++(1.5,0) coordinate(Out1) to[short, -*] ++(2,0) coordinate(Out2) to[short, -o] ++(0.5,0) node[right]{$\text{V}_{\text{out}}$}
      				(OA.-) node[below] {$2$} to[short,] ++(0,+1) to[short,] ++(2.75,0) to[short, -*] ++(0,-1.5);
      				\draw (-4,-2) node[ground]{} to[sinusoidal voltage source] ++ (0,+2)
      				to[R=$4.7 \text{ k}\Omega$, -*] ++(2,0) coordinate(In1)
      				to[R=$4.7 \text{ k}\Omega$] ++(0,-2) node[ground]{}
      				(Out2) to[oscope=$\text{CH2}$, -] ++(0,-2.5) node[ground]{}
      				(Out1) to[R=$10 \text{ k}\Omega$] ++(0,-2.5) node[ground]{}
      				(NonInv) to[oscope=$\text{CH1}$, -] ++(0,-2) node[ground]{}
      				(NonInv) to[short,] ++(-2,0)
      				(NonInv) node[above] {$\text{V}_{\text{in}}$}
      				(In1) to[short,] ++(1.65,0) to[R=$1 \text{ k}\Omega$] ++(0,-2) node[ground]{};
      			\end{circuitikz}
  \caption{Buffer Circuit \#3 with the 10~k$\Omega$ load resistor and a 1~k$\Omega$ resistor in parallel with the lower 4.7~k$\Omega$ voltage divider resistor.}
\end{figure}

\clearpage

\noindent\textbf{Label:} \texttt{fig:lm6134-pins2}\\[0.5em]
\begin{figure}[H]
  \centering
  \begin{circuitikz}[american, scale=1]
      			% --- AMPLIFIER A ---
      			\draw (0,0) coordinate(NonInvA) to[short, o-] ++ (1,0)
      			node[above] {$3$}
      			node[op amp, noinv input down, anchor=+](OA){OA}
      			(OA.out) node[above] {$1$} to[short, -o] ++(0.5,0)
      			node[right]{$\text{OUT A}$}
      			(OA.-) node[above] {$2$} to[short, -o] ++(-1,0)
      			coordinate(InvA)
      
      			% Power Connections for Amp A (Shared for IC)
      			(OA.up) coordinate(posP)
      			to[short, -*] ++(0,1.25) coordinate(one)
      			to[short, -o] ++(0,0.5)
      			node[above] {$8$} node[left] {$\text{V}^+$}
      			(one) to[curved capacitor=$0.1\,\mu\text{F}$]
      			++(2,0) node[ground]{}
      
      			(OA.down) coordinate(posN)
      			to[short, -*] ++(0,-1.25) coordinate(two)
      			to[short, -o] ++(0,-0.5)
      			node[below] {$4$} node[left] {$\text{V}^-$}
      			(two) to[curved capacitor=$0.1\,\mu\text{F}$]
      			++(2,0) node[ground]{}
      
      			(NonInvA) node[left] {$\text{+IN A}$}
      			(InvA) node[left] {$\text{-IN A}$};
      
      			% --- AMPLIFIER B ---
      			\draw (7,0) coordinate(NonInvB) to[short, o-] ++ (1,0)
      			node[above] {$5$}
      			node[op amp, noinv input down, anchor=+](OB){OB}
      			(OB.out) node[above] {$7$} to[short, -o] ++(0.5,0)
      			node[right]{$\text{OUT B}$}
      			(OB.-) node[above] {$6$} to[short, -o] ++(-1,0)
      			coordinate(InvB)
      
      			% Power Connections for Amp B (linked to same pins)
      			(OB.up) to[short, -*] ++(0,1.25)
      			(OB.down) to[short, -*] ++(0,-1.25)
      
      			(NonInvB) node[left] {$\text{+IN B}$}
      			(InvB) node[left] {$\text{-IN B}$};
      		\end{circuitikz}
  \caption{LM662 connection diagram. The $0.1\,\mu$F decoupling capacitors are required on every build. Because pins~4 and~8 are shared across both op-amps, one capacitor pair placed adjacent to the supply pins protects both amplifiers simultaneously.}
\end{figure}

\end{document}
